% SHARD-ID: AQM-89-SCHEME-INCIDENCE-CODE-SOURCES
% SHARD-TITLE: Scheme-Incidence Sources for Code Complexes
% SHARD-SUMMARY: Separates bare finite spectra from schemes equipped with incidence or group-law data.
% SHARD-SUMMARY: Distinguishes checked Weyl-label selections from sourced targets, proposals, and unreconstructed naturality claims.
% SHARD-SUMMARY: Explains why the constant finite group-scheme toric route is structured extra data, not a consequence of \(\Spec R\) alone.
% SHARD-KEYWORDS: scheme incidence, group scheme, toric code, Steane code, Reed-Muller, Shor code, Weyl-Heisenberg constraints

\section{Scheme-Incidence Sources for Code Complexes}
\label{sec:scheme-incidence-code-sources}

\subsection{Status and Source Anchors}
This shard answers whether rings or schemes can give the complexes proposed
after AQM-88.  The answer is conditional: a bare finite ring gives finite
residue Weyl sites, while a code complex needs extra selected functions,
incidence relations, or group-law data whose boundary identity selects a
commuting Weyl-Heisenberg subgroup.  The CSS bridge is AQM-08.  It proves that
a finite chain complex
\[
  C_2\xrightarrow{\partial_2}C_1\xrightarrow{\partial_1}C_0,
  \qquad \partial_1\partial_2=0,
\]
selects the Weyl-label subgroup
\[
  L=(\operatorname{im}\delta^0\oplus0)+(0\oplus\operatorname{im}\partial_2)
  \subset C^1\oplus C_1 .
\]
AQM-10 fixes the Steane binary presentation; AQM-05 and AQM-06 validate the
periodic square toric complex for \(k=4\); AQM-40 proves
\(\Spec(k^n)\); AQM-88 requires final selected/rejected Weyl-Heisenberg
elements.  The convention is \texttt{CONVENTIONS.md} \texttt{(ax)}.  The only
group-scheme material used here is the explicit constant finite case
\(G_k=\Spec(k^G)\).  Stacks supplies affine-scheme anti-equivalence in
\path{references/algebraic_geometry/stacks_project_schemes.tex}, lines
841--963, and finite products/fibre products in lines 986--1004.  No general
nonconstant group-scheme theory is imported.

\subsection{Lamport Dependency Skeleton}

\[
\begin{array}{rcl}
D1&:&\Spec R\text{ gives residue sites, not a complex.}\\
D2&:&X_i=\Spec(k^{S_i}),\ I_{10},I_{21}\text{ define }\partial_1,\partial_2.\\
L1&:&\partial_1\partial_2=0\Rightarrow
L=(\operatorname{im}\delta^0\oplus0)+(0\oplus\operatorname{im}\partial_2)
\text{ is isotropic.}\\
T1&:&\text{Steane, Reed-Muller, Shor, toric are separated below by evidence status.}
\end{array}
\]

\subsection{Why a Bare Finite Ring Is Not Enough}

Finite rings have finitely many ideals, hence satisfy DCC and are Artinian.
Stacks \path{references/algebraic_geometry/stacks_project_algebra.tex}, lines 12727--12805 and 14364--14444, records finite maximal ideals, maximal primes, and finite discrete spectrum after local decomposition.  AQM-40 gives the concrete reduced case
\[
  \Spec(k^n)=\{\mathfrak m_1,\ldots,\mathfrak m_n\},
  \qquad \kappa(\mathfrak m_i)=k.
\]
This gives \(E=\bigoplus_i k^2\), but it does not supply edges, faces, or a
boundary map.  Stabilizer constraints are extra isotropic label choices.  Thus
the correct outside-in target is not \(\Spec R\leadsto C_2\to C_1\to C_0\)
without further data, but
\[
  \text{finite schemes plus named incidence/group structure}
  \leadsto C_2\to C_1\to C_0
  \leadsto \pi^{-1}(L)\subset\operatorname{WH}(E).
\]

\subsection{Steane: Punctured Boolean Three-Space}

Let \(X_7=\mathbb F_2^3\setminus\{0\}\), ordered as
\[
  111,\ 110,\ 101,\ 100,\ 011,\ 010,\ 001 .
\]
The finite scheme is \(X_{\rm St}=\Spec(\mathbb F_2^{X_7})\), equivalently
Boolean affine three-space with the origin idempotent removed:
\[
  \mathbb F_2[x_1,x_2,x_3]/(x_i^2+x_i,\ (1+x_1)(1+x_2)(1+x_3)).
\]

Evaluate the three coordinate functions on the displayed order.  The rows are
\[
\begin{aligned}
  x_1&=1111000,\\
  x_2&=1100110,\\
  x_3&=1010101.
\end{aligned}
\]
These are exactly the Steane rows \(g_1,g_2,g_3\) in AQM-10.  Put
\(D=\langle x_1,x_2,x_3\rangle\subset\mathbb F_2^{X_7}\).  Each row has weight
\(4\), and any two distinct rows overlap in \(2\) sites, so
\[
  D\subset D^\perp .
\]
Therefore the selected Weyl-label subgroup
\[
  L_{\rm St}=(D\oplus0)+(0\oplus D)
\]
is isotropic.  AQM-10 then identifies \(L_{\rm St}^\perp/L_{\rm St}\) with one
binary symplectic logical pair.

The nontrivial selection principle is the choice of homogeneous degree-one
functions \(\langle x_1,x_2,x_3\rangle\), not all affine functions.  Including
the constant function would change the row space.  Hence the ring/scheme
source is promising but not yet a dynamical derivation of the selection rule.

\subsection{Reed-Muller: Boolean Degree Filtration}

Let
\[
  B_m=\mathbb F_2[x_1,\ldots,x_m]/(x_i^2+x_i),\qquad X_m=\Spec B_m .
\]
Then \(X_m\) is the finite Boolean \(m\)-cube and \(B_m\) is the full function
ring on it.  For \(S\subset\{1,\ldots,m\}\), write
\[
  x_S=\prod_{i\in S}x_i .
\]
The squarefree monomials form the function basis.  Define the degree-filtered
subspace
\[
  RM(r,m)=\langle x_S:\ |S|\le r\rangle\subset\mathbb F_2^{\mathbb F_2^m}.
\]
This shard uses this displayed definition as the local Reed-Muller convention.

Use the pointwise binary pairing
\[
  \langle f,g\rangle=\sum_{a\in\mathbb F_2^m} f(a)g(a).
\]
For monomials,
\[
  \langle x_S,x_T\rangle
  =
  \sum_a x_{S\cup T}(a)
  =
  \begin{cases}
  1,&S\cup T=\{1,\ldots,m\},\\
  0,&S\cup T\ne\{1,\ldots,m\}.
  \end{cases}
\]
Indeed, \(x_{S\cup T}\) is \(1\) on \(2^{m-|S\cup T|}\) points, and this number
is odd exactly when \(|S\cup T|=m\).

Hence \(RM(r,m)\) is orthogonal to \(RM(s,m)\) whenever \(r+s<m\).  For \(0\le r<m\),
\[
  RM(r,m)^\perp=RM(m-r-1,m),
\]
because the inclusion follows from the monomial calculation and the dimensions
sum to \(2^m\).

Therefore every pair \(r_X+r_Z<m\) gives a checked CSS Weyl-label selection
\[
  L_{r_X,r_Z}
  =
  (RM(r_X,m)\oplus0)+(0\oplus RM(r_Z,m)).
\]
The selected Weyl-Heisenberg elements are \(\pi^{-1}(L_{r_X,r_Z})\).  Any
stronger Reed-Muller-code claim, for example a particular distance or
transversal-gate property, still needs its own local source or run.

\subsection{Shor: A Sourced Incidence Target, Not Yet a Natural Scheme}

Gottesman's local source lists Shor's nine-qubit stabilizer table in
\path{references/symplectic_qecc/gottesman_9705052_source/Thesis.tex}, lines
920--935.  In binary CSS row notation, take \(S_1=\{1,\ldots,9\}\) and
\[
\begin{gathered}
H_X=\begin{pmatrix}111111000\\111000111\end{pmatrix},\\
H_Z=\begin{pmatrix}
110000000\\101000000\\000110000\\000101000\\000000110\\000000101
\end{pmatrix}.
\end{gathered}
\]
Each \(Z\)-row is supported inside one block of three.  Each \(X\)-row is
constant on any block it touches.  Thus every mixed dot product is either \(0\)
or \(2\), hence zero in \(\mathbb F_2\):
\[
  H_XH_Z^T=0.
\]

This gives a finite incidence chain complex by setting
\[
  C_0=\mathbb F_2^2,\quad C_1=\mathbb F_2^9,\quad C_2=\mathbb F_2^6,
  \quad \partial_1=H_X,\quad \partial_2=H_Z^T .
\]
The selected Weyl labels are
\[
  L_{\rm Shor}
  =
  (\operatorname{row}H_X\oplus0)+(0\oplus\operatorname{row}H_Z).
\]
The rows have ranks \(2\) and \(6\), so \(L_{\rm Shor}\) has rank \(8\) inside
\(\mathbb F_2^9\oplus\mathbb F_2^9\).

This is a valid scheme-incidence object in convention \texttt{(ax)}: take
\(\Spec(\mathbb F_2^{S_i})\) for the check sets and encode the ones of
\(H_X,H_Z\) as incidence relations.  The open question is naturality.
Bombin--Martin--Delgado support the homological-family statement locally in \path{references/toric_code/bombin_martin_delgado_0605094_source/HomologicalErrorCorrection6.tex}, lines 320--321 and 2488--2493; this shard still does not reconstruct their natural surface \(2\)-complex.

\subsection{Toric: Constant Group Scheme and Cayley Squares}

Let \(G=C_k\times C_k\) be the finite abelian group written additively, with
generators \(a=(1,0)\) and \(b=(0,1)\).  Over \(\mathbb F_2\), the constant
finite group scheme used here is \(G_{\mathbb F_2}=\Spec(\mathbb F_2^G)\).  As
a bare scheme this is only the product-field spectrum with \(k^2\) points.  The
group structure is extra morphism data.  In coordinate functions
\(\delta_g\in\mathbb F_2^G\), multiplication is the affine-scheme morphism with
pullback
\[
  m^\#(\delta_h)=\sum_{u+v=h}\delta_u\otimes\delta_v .
\]
Unit and inverse are defined similarly.  Thus the constant group scheme is a
scheme, but not merely \(\Spec(\mathbb F_2^G)\) with its extra structure
forgotten.

Build the periodic square incidence complex:
\[
  S_0=G,\qquad S_1=G\times\{a,b\},\qquad S_2=G .
\]
For \(g\in G\), use edges \((g,a):g\to g+a\) and \((g,b):g\to g+b\).  Over
\(\mathbb F_2\) signs disappear, so set
\[
  \partial_1(g,a)=g+(g+a),\qquad
  \partial_1(g,b)=g+(g+b),
\]
and
\[
  \partial_2(g)=(g,a)+(g+a,b)+(g+b,a)+(g,b).
\]
Then
\[
\begin{aligned}
\partial_1\partial_2(g)
&=g+(g+a)+(g+a)+(g+a+b)\\
&\quad +(g+b)+(g+a+b)+g+(g+b)\\
&=0 .
\end{aligned}
\]
The cancellation is exactly the commuting-square identity \(g+a+b=g+b+a\).

Put qubit Weyl labels on \(S_1\).  The toric selected label subgroup is
\[
  L_{\rm tor}(G,a,b)
  =
  (\operatorname{im}\delta^0\oplus0)+(0\oplus\operatorname{im}\partial_2)
  \subset \mathbb F_2^{S_1}\oplus\mathbb F_2^{S_1}.
\]
The selected Weyl-Heisenberg elements are \(\pi^{-1}(L_{\rm tor})\).  For
\(k=4\), this gives \(|S_0|=16\), \(|S_1|=32\), and \(|S_2|=16\), matching the
AQM-05/AQM-06 validated toric instance.  The Haah/Laurent-polynomial layer is
deferred to AQM-91 and is not used by this finite Cayley-incidence precursor.

\paragraph{Is this cheating?}
It is not cheating if the input is allowed to be a scheme with named group law
and named generators.  It is cheating if one says that the bare ring
\(\mathbb F_2^G\), or the bare scheme \(\Spec(\mathbb F_2^G)\), canonically
produces the toric code.  The bare scheme only knows \(k^2\) points.  The toric
complex uses three extra choices:
\[
  \text{group law on }G,\qquad
  \text{commuting generators }a,b,\qquad
  \text{Cayley-square face convention}.
\]
These choices are substantial, but they are structured.  They replace an
arbitrary \(32\)-qubit CSS matrix by a translation-invariant incidence complex
whose CSS commutation follows from \(a+b=b+a\).

\subsection{Conclusion}

The evidence separates the four routes: Steane is punctured Boolean
three-space plus a degree-one selection rule; Reed-Muller is the Boolean cube
with a degree filtration whose orthogonality is checked here, while distance
and gate claims still need sources or runs; Shor has the displayed incidence
target and sourced BMD family locator, but the natural surface complex is not
reconstructed here; toric is the constant group scheme plus the finite
Cayley-square complex, with the arithmetic status of the group/incidence input
still to justify.  In all four cases the final object is a selected
Weyl-label subgroup \(L\), hence \(\pi^{-1}(L)\subset\operatorname{WH}(E)\).
